\begin{proof}
\begin{enumerate}[label=\arabic*.,leftmargin=*]
\item Use $P_E=P_{\rm mark}$ from \cref{inv:B1.0}. In Engel's marking,
the four columns of the change to Viaclovsky's basis are
\[
 Q_E=\begin{pmatrix}-2&1&0&-2\\-1&0&0&0\\3&-1&1&0\\0&0&0&1\end{pmatrix}.
\]
Direct integral multiplication gives $P_V=Q_E^{-1}P_E$ and
$V_0P_V=P_VA_1$, $V_1P_V=P_VA_2$, where $V_0,V_1$ are the
matrices in that paper's Lemma 2.1.  Multiplication by the three
specified invariant vectors gives the displayed identities.
Inverting a clockwise affine generator replaces its invariant
translation $w_j/m_j$ by $-w_j/m_j$.  Consequently the
counterclockwise translation vector is $-P_Vw_j$, which gives the
two formulas, including the term $N_1$ in the second coefficient.  Viaclovsky's reference section extends through
the smooth locus of the cusp; its comparison therefore has zero
class in the quotient lattice of toric directions.  This is precisely
$l_0=0$.  Substitution into the definition of $p$ proves its formula.

\end{enumerate}
\end{proof}